% Build from the repository root:
% latexmk -pdf -interaction=nonstopmode -halt-on-error \
%   -outdir=/tmp/fulton-harris-description-build docs/fulton-harris-pilot-description.tex
% cp /tmp/fulton-harris-description-build/fulton-harris-pilot-description.pdf docs/
\documentclass[11pt,a4paper]{article}
\usepackage[T1]{fontenc}
\usepackage{lmodern}
\usepackage[margin=23mm,headheight=15pt]{geometry}
\usepackage{amsmath,amssymb,mathtools}
\usepackage{microtype}
\usepackage[table]{xcolor}
\usepackage{booktabs,tabularx,array}
\usepackage{enumitem}
\usepackage{fancyhdr}
\usepackage{xurl}
\usepackage[colorlinks=true,linkcolor=teal!60!black,urlcolor=teal!60!black]{hyperref}
\definecolor{ink}{HTML}{183947}
\definecolor{muted}{HTML}{53636A}
\pagestyle{fancy}
\fancyhf{}
\fancyhead[L]{\small\textcolor{ink}{Fulton--Harris pilot: mathematical description}}
\fancyhead[R]{\small\textcolor{muted}{20 October 2026}}
\fancyfoot[C]{\small\thepage}
\setlength{\parindent}{0pt}
\setlength{\parskip}{6pt plus 1pt}
\setlength{\emergencystretch}{2em}
\setlist{leftmargin=*,itemsep=3pt,topsep=4pt,parsep=0pt}
\renewcommand{\arraystretch}{1.2}
\newcolumntype{Y}{>{\raggedright\arraybackslash}X}
\newcommand{\Q}{\mathbb Q}
\newcommand{\C}{\mathbb C}
\newcommand{\Hom}{\operatorname{Hom}}
\newcommand{\End}{\operatorname{End}}
\newcommand{\tr}{\operatorname{tr}}
\newcommand{\id}{\operatorname{id}}
\newcommand{\code}[1]{\texttt{\detokenize{#1}}}
\newcommand{\source}[1]{\par{\small\raggedright\color{muted}Implementation: #1.\par}}
\hypersetup{pdftitle={Fulton--Harris pilot: a mathematical description},
  pdfsubject={Constructive representations, split submodules, characters, and S3},
  pdfauthor={Arend formalization project}}

\begin{document}
\thispagestyle{empty}
{\color{ink}\LARGE Fulton--Harris pilot\par}
{\Large A mathematical description of the formalization\par}
{\color{muted}20 October 2026\par}

The pilot develops a constructive version of selected results from the first
two lectures of Fulton--Harris, inside \code{arend-lib}. Its worked example is
the representation theory of $S_3$: the trivial, sign, standard and permutation
representations, an explicit decomposition, and character computations.
This document explains the current mathematics; the detailed implementation
plan is in \href{fulton-harris-pilot-plan.pdf}{\code{fulton-harris-pilot-plan.pdf}}.

\textbf{Representations and their maps.}
A representation consists of an $R$-module $V$ and a linear action of a group
$G$. Writing the action as $gv$, its laws include
\[
  (gh)v=g(hv),\qquad 1v=v,\qquad
  g(v+w)=gv+gw,\qquad g(av)=a(gv).
\]
An intertwiner $f:V\to W$ is a linear map satisfying $f(gv)=g f(v)$.
The usual constructions are present: direct sums, tensor products, duals,
and the action on linear maps
\[
  (g\cdot f)(v)=g f(g^{-1}v),\qquad
  V^*=\Hom_R(V,R),
\]
where $R$ has the trivial action. The underlying objects are modules.
Many constructions work over commutative rings; stronger assumptions are
attached to the particular theorems that need them.

\textbf{One hierarchy of subrepresentations.}
The existing class \code{SubLRepres} describes an invariant submodule by a
membership predicate $P(v)$ and its closure laws. Its underlying representation
is the canonical subtype
\[
                         W=\{v\in V\mid P(v)\},
\]
with inclusion $j(v,p)=v$. The new descendant \code{SplitSubLRepres} adds a
linear retraction $r:V\to W$ satisfying $r\circ j=\id_W$. It inherits the
membership predicate and inclusion; it does not introduce a different
subrepresentation. The retraction need not respect the group action.
The map $p=j\circ r$ is a projection, since $p^2=p$.

The module classes follow the same pattern. In code, \code{SubLModule} and
\code{SubLRepres} coerce to the module or representation on their members.
Thus \code{x : W} denotes a member of $W$, while \code{W v} is the proposition
that an ambient vector $v$ belongs to $W$. The ambient object is \code{W.S}.
The split and based descendants inherit these coercions. The based descendants
additionally record finite free coordinates for $W$.

\textbf{What is assumed about $W$.}
To test irreducibility, we consider an invariant submodule $W\subseteq V$
and a supplied linear map $r:V\to W$ that fixes $W$: $r(w)=w$ for every
$w\in W$. This is the splitting stored by \code{SplitSubLRepres}.
No basis of $W$ is assumed. A finite basis of $V$ is assumed to exist.

A basis of $W$ can be supplied as additional data for coordinate computations.
The based classes record this extra data. It is not needed for the
irreducibility condition.

\textbf{The working irreducibility notion.}
The class \code{FinIrreducibleRepres} extends \code{FinLinRepres}, which
inherits finite freeness from the existing \code{FinModule} class. Thus a finite
basis of $V$ is assumed to exist; no particular basis is chosen. Its
\code{nontrivial} field asserts $\exists v\in V,\ v\ne0$. Together these
properties imply that every basis is nonempty, as proved by
\code{basis-nonempty}.

Its \code{irreducible} field says that for every split subrepresentation
$W\subseteq V$ with its retraction, either every element of $W$
is zero, or $j$ is surjective. Thus every invariant submodule with a linear
retraction is zero or whole. This is the pilot's finite free notion of
irreducibility; a general notion for arbitrary modules is outside its scope.

\source{\code{SubLRepres}, \code{SplitSubLRepres}, \code{BasedSplitSubLRepres},
\code{FinLinRepres}, \code{FinIrreducibleRepres.splitMono-irreducible}}

\newpage
\textbf{Why a splitting is part of the hypothesis.}
Finite dimension of the ambient space does not, by itself, constructively
supply a complement to a submodule described by an arbitrary predicate.
For any proposition $P$, consider the submodule of the trivial representation
on a nonzero ring
\[
                   W_P=\{x\in R\mid P\lor x=0\}.
\]
The zero alternative ensures that this is a submodule even when $P$ is false.
If $W_P$ is zero, then $P$ is false, because $P$ would put $1$ in $W_P$.
If $W_P$ is whole, membership of $1$ implies $P$. Consequently, a theorem
asserting that every such predicate submodule is zero or whole would decide
every proposition. This implication is formalized for the old definition,
now named \code{IsStronglyIrreducible}.

The split formulation only applies when the retraction is supplied or proved
to exist. No general procedure producing a retraction for every $W_P$ is
assumed. The distinction concerns the information available about a
submodule, including in a finite-dimensional ambient space.

\textbf{Maschke's averaging argument.}
Let $G$ be finite and suppose an inverse of the scalar $|G|\cdot1_R$ is given.
Starting with the linear retraction above, define
\[
                 \bar r(v)=\frac1{|G|}\sum_{g\in G}g\,r(g^{-1}v).
\]
The implementation proves that $\bar r$ is an intertwiner and that
$\bar r\circ j=\id_W$. Thus a splitting of modules becomes a splitting of
representations. Equivalently, averaging $j\circ r$ gives an equivariant
projection $\bar p$ onto $j(W)$.

The expected decomposition is
\[
 v=\bar p(v)+(v-\bar p(v)),\qquad
 \bar p(v-\bar p(v))=0.
\]
The averaging theorem works over a commutative ring with the supplied inverse
and needs no finite basis. The implemented theorem constructing the
complement and an isomorphism $V\cong W\oplus C$ additionally assumes a finite
basis of $V$ and the library's \code{SmithDomain} linear-algebra structure.
This includes discrete fields; it is not an assertion of such a basis
construction over arbitrary constructive $\C$.

\textbf{Proving irreducibility through scalar endomorphisms.}
Suppose every intertwiner $V\to V$ is multiplication by a scalar.
For a split invariant submodule, the averaged projection therefore has the form
\[
               \bar p=c\,\id_V,\qquad c^2=c.
\]
Over a commutative local ring, the scalar $c$ is zero or one. If it is zero,
the submodule is zero; if it is one, the inclusion is surjective. This proves
the pilot's irreducibility statement. The argument decides only a scalar
already known to be idempotent; it needs no test of whether an arbitrary
scalar is zero. This is the route used for the standard representation of
$S_3$ over constructive $\C$.

\source{\code{IsStronglyIrreducible.excludedMiddle}, \code{maschke-split},
\code{maschke-complement}, \code{irreducible-from-commutant}}

\newpage
\textbf{Schur's lemma and its constructive boundary.}
Over a discrete field $K$ (a field with decidable equality), the implementation
proves that an intertwiner between irreducible representations is zero or an
isomorphism. The theorem takes irreducible representations directly, rather than a
representation and a separate proof of an irreducibility predicate.
Its proof constructs split presentations of the kernel and image.
Irreducibility then says that the image is zero or whole, and that the kernel
is zero or whole. The resulting cases give the conclusion.

The unrestricted zero-or-isomorphism conclusion is not asserted over
constructive $\C$. Already on the one-dimensional trivial representation,
multiplication by any scalar $c$ is an intertwiner. Such a dichotomy would
tell us whether $c=0$ or $c$ is invertible, and hence decide scalar equality.
The source records this obstruction as \code{schur-decideEq}. It is an
implication from the proposed theorem, not a proof of its negation.

For complex character theory the available route keeps the useful
hypotheses explicit: all endomorphisms are scalar, or all intertwiners
between two specified representations vanish. Deriving these hypotheses
from irreducibility in full generality remains a separate problem.

\textbf{What information permits division.}
The implementation distinguishes three sorts of evidence:
\begin{itemize}
  \item an equality $a=0$;
  \item its negation $\neg(a=0)$;
  \item positive invertibility evidence: a scalar $b$ with $ab=ba=1$.
\end{itemize}
Over a discrete field the relevant zero tests and inverse construction are
available. Over constructive $\C$, a negated equality is not silently turned
into an inverse or into positive apartness evidence. When a computation
divides, its hypotheses must justify that division.

For example, a pivot test in elimination asks whether a matrix entry $a$ is
zero, then uses an inverse in the nonzero branch. This is stronger than merely
performing additions and multiplications. Such elimination belongs to the
discrete-field part of the pilot. Averaging instead takes its required inverse
as an explicit input.

\textbf{The passage from rational to complex coefficients.}
The pilot constructs the ring homomorphism $\Q\to\C$ and uses it to transport
inverses of positive integers. In particular, the complex inverses of $3$
and $6$ come from their rational inverses. No search for an arbitrary
complex inverse is needed.

The formulas for the $S_3$ representations, their decomposition and their
characters are proved with general coefficients and instantiated at $\C$.
The standard representation's complex irreducibility has the scalar-projection
proof described above. It is not inferred merely from irreducibility over
$\Q$: irreducibility need not survive a field extension in general.
General scalar-extension machinery is not part of this pilot.

\textbf{Rank one.}
A representation free of rank one is proved irreducible in the split sense
when the coefficient ring has only the idempotents $0$ and $1$. This proof
does not need averaging or an inverse of $|G|$. The condition matters over
rings: a product ring has proper direct summands even as a rank-one module
over itself.

\source{\code{SchurLemma}, \code{schur-decideEq}, \code{rank-one-irreducible};
\code{Arith/Complex/RatHom.ard}}

\newpage
\textbf{Characters and their basic formulas.}
For a finite free representation over a commutative ring, define
\[
                      \chi_V(g)=\tr(\rho_V(g)).
\]
The trace is proved independent of the chosen basis. The resulting character
is constant on conjugacy classes, and $\chi_V(1)$ is the dimension interpreted
as a scalar in the coefficient ring. The implementation proves
\begin{align*}
 \chi_{V\oplus W}(g)&=\chi_V(g)+\chi_W(g),\\
 \chi_{V\otimes W}(g)&=\chi_V(g)\chi_W(g),\\
 \chi_{V^*}(g)&=\chi_V(g^{-1}),\\
 \chi_{\Hom_R(V,W)}(g)&=\chi_V(g^{-1})\chi_W(g).
\end{align*}

\textbf{A character without a basis of the submodule.}
If $j:W\to V$ has a linear retraction $r$, and $V$ is finite free, define
\[
           \chi_W(g)=\tr_V\bigl(j\circ\rho_W(g)\circ r\bigr).
\]
This uses a trace on the ambient module. It does not require a basis of $W$.
The formalization proves independence of the choice of retraction, agreement
with the ordinary character if a basis of $W$ is available, agreement with the chosen coordinate
presentation (\code{splitCharacter-coordinates}), and additivity
across a supplied decomposition into split subrepresentations. These are the
implemented results in \code{SplitCharacter.ard}.

\textbf{Averaging, invariants and orthogonality.}
The map
\[
                 P(v)=\frac1{|G|}\sum_{g\in G}gv
\]
is proved to project onto the invariant vectors. Its trace is
$|G|^{-1}\sum_g\chi_V(g)$. If those invariant vectors have a supplied finite
free split presentation of rank $d$, the trace is also $d\cdot1_R$.
The existence of such a presentation is proved over discrete fields.

The character pairing is defined algebraically by
\[
       \langle\chi,\psi\rangle
       =\frac1{|G|}\sum_{g\in G}\chi(g^{-1})\psi(g).
\]
Applying the averaging trace formula to $\Hom_R(V,W)$ gives the route to
orthogonality. Vanishing of all intertwiners implies pairing zero; a scalar
endomorphism space gives self-pairing one for a nonzero finite free
representation. The library also deduces the non-isomorphic irreducible case
over discrete fields using Schur's lemma. These general orthogonality results
depend on Exercise~1.2, which is deliberately left as a hole.

For a supplied finite family of pairwise non-isomorphic irreducibles with
scalar endomorphisms, the implementation uses characters to prove uniqueness
of multiplicities in two given isomorphic direct sums. This theorem assumes
a discrete field and injectivity of the natural-number map into it
(characteristic zero). It inherits the exercise dependency. It does not
construct a decomposition of an arbitrary representation.

\source{\code{Characters.ard}, \code{CharacterFormulas.ard}, \code{Tensor.ard},
\code{SplitCharacter.ard}, \code{Invariants.ard}, \code{Orthogonality.ard},
\code{Multiplicities.ard}}

\newpage
\textbf{The explicit $S_3$ example.}
The permutation representation is the action on $R^3$ by permuting
coordinates. Its standard subrepresentation is the sum-zero plane
\[
                 W=\{(x,y,z)\in R^3\mid x+y+z=0\}.
\]
It has explicit coordinates and a retraction
\[
       j(a,b)=(a,b-a,-b),\qquad r(x,y,z)=(x,-z).
\]
Here $rj=\id$, and $jr$ is the identity on the sum-zero plane. This
presentation uses no division. When $3$ is invertible, put
$m=(x+y+z)/3$. The equivariant decomposition is
\[
 (x,y,z)=m(1,1,1)+(x-m,y-m,z-m),\qquad
 R^3\cong R_{\mathrm{triv}}\oplus W.
\]
Both directions of the isomorphism and their inverse laws are implemented.
An explicit calculation shows that every equivariant endomorphism of $W$
is scalar. Averaging then proves irreducibility over a commutative local ring
with $6$ invertible, including $\Q$ and constructive $\C$.

The computed character values are:
\begin{center}
\begin{tabular}{@{}lrrr@{}}
\toprule
Representation & Identity & Transposition & Three-cycle\\
\midrule
Trivial     & $1$ & $1$  & $1$\\
Sign        & $1$ & $-1$ & $1$\\
Standard    & $2$ & $0$  & $-1$\\
Permutation & $3$ & $1$  & $0$\\
\bottomrule
\end{tabular}
\end{center}
The six permutations are explicitly enumerated, and all nine pairings among
the trivial, sign and standard characters are proved: one on the diagonal,
zero off it. These direct calculations are independent of the exercise holes
and apply over constructive $\C$ using the supplied inverse of $6$.
The pilot does not yet prove that these exhaust all irreducible
representations of $S_3$.

\textbf{The intentional holes.}
Two exercises remain stated with unfilled proofs, as requested:
\begin{itemize}
 \item Exercise~1.2: invariant elements of $\Hom_R(V,W)$ are exactly
       intertwiners. Its goal is \code{FH_L01_Ex1_2}.
 \item Exercise~2.5: the character of a permutation representation counts
       fixed points. Its goal is \code{FH_L02_Ex2_5}.
\end{itemize}
Any result using these holes is conditional on the corresponding exercise.
The direct $S_3$ computations and its irreducibility proof do not use them.

\textbf{The present boundary.}
The pilot supplies constructive representations, split irreducibility,
averaging, characters and the explicit $S_3$ results under their stated
hypotheses. General scalar extension, a general decomposition algorithm,
completeness of the $S_3$ list, Hermitian forms, and the general complex
consequences of Schur's lemma from irreducibility alone remain outside these
implemented results.

\source{\code{Symmetric3.ard}, \code{Symmetric3Characters.ard};
\code{test/Algebra/Representation/S3Test.ard}}
\end{document}
